\documentclass[11pt]{article}
\usepackage[margin=1in]{geometry}
\usepackage{hyperref}
\usepackage{enumitem}
\usepackage{booktabs}

\title{Loop Engine: Verifiable Closed-Loop Runtime for Research Agents}
\author{Research Network Loop Lab}
\date{June 2026}

\begin{document}
\maketitle

\begin{abstract}
Agentic research systems increasingly operate as loops: observe a workspace, plan a step,
execute tools, critique the result, and publish a new artifact. Most systems treat that loop
as transient orchestration hidden inside a chat transcript or private runtime log. Loop Engine
turns the loop into a verifiable research object. Each cycle emits a typed state transition,
an evidence ledger, artifact hashes, reviewer checks, and graph edges that can be stored on
Walrus and anchored through Sui. The result is a runtime pattern where autonomous research is
not merely answered, but made replayable, forkable, citeable, and economically accountable.
\end{abstract}

\section{Problem}

The unit of AI-assisted research is no longer a single paper or a single answer. A useful
agent may browse sources, inspect code, call tools, run tests, write artifacts, consult other
agents, revise the plan, and eventually publish a package. The value is distributed across the
loop, but the public record usually preserves only the final response.

This creates a verification gap. Readers cannot tell which external evidence was used, which
tool calls changed the workspace, which failed paths were rejected, or why the final artifact
should be trusted. Downstream agents face the same problem at higher speed: they can reuse an
output, but not the reasoning workflow that produced it.

\section{Loop Engine}

Loop Engine models an autonomous research run as a sequence of state transitions:

\begin{enumerate}[leftmargin=*]
  \item \textbf{Observe}: bind the current workspace, graph context, chain state, and external evidence.
  \item \textbf{Plan}: declare the next step and the acceptance condition before execution.
  \item \textbf{Execute}: run tools, edits, simulations, browser actions, or chain transactions.
  \item \textbf{Critique}: attach checks, reviewer objections, failed hypotheses, and residual risk.
  \item \textbf{Publish}: commit accepted deltas into a Research Asset manifest, Walrus snapshot, and Sui registry event.
\end{enumerate}

The engine does not require every token of private reasoning to be public. It requires that
public claims be backed by enough typed evidence for a reader or another agent to replay the
operational path.

\section{Protocol Mapping}

The Research Network stack gives the loop a durable surface. Git stores the working tree.
Walrus stores release archives and loop evidence bundles. Sui records registry events,
object identifiers, manifest hashes, and settlement hooks. The static web surface renders the
human view, while the SDK and skill package expose the agent-readable view.

\begin{center}
\begin{tabular}{lll}
\toprule
Loop phase & Protocol artifact & Verification handle \\
\midrule
Observe & evidence bundle & source hash, timestamp, URI \\
Plan & loop intent & acceptance condition hash \\
Execute & artifact delta & file checksum, tool receipt \\
Critique & review record & test result, reviewer signature \\
Publish & Research Asset release & Walrus blob, Sui object, manifest hash \\
\bottomrule
\end{tabular}
\end{center}

\section{Why It Matters}

Loop Engine changes the default public artifact from ``what the agent said'' to ``what the
agent changed and why that change survived review.'' This matters for research because reuse is
becoming agentic. A future agent should be able to install a skill, replay a workflow, inspect
which claims were grounded, and fork the asset without losing provenance.

For marketplaces and tokenized research networks, the same loop record also becomes an
accounting primitive. If reuse, citation, or delegated work carries economic value, the loop
record explains which steps and artifacts earned attribution.

\section{Evaluation}

A Loop Engine implementation should be evaluated on four properties:

\begin{enumerate}[leftmargin=*]
  \item \textbf{Replayability}: another agent can reconstruct the public sequence of state transitions.
  \item \textbf{Grounding}: each published claim points to evidence or an explicit assumption.
  \item \textbf{Forkability}: downstream work can inherit graph edges and license terms without manual archaeology.
  \item \textbf{Settlement readiness}: reusable skills, reports, and workflow outputs expose enough identifiers for later economic routing.
\end{enumerate}

\section{Conclusion}

Loop Engine is a runtime pattern for making autonomous research legible. It does not replace
papers; it wraps them in the loop evidence needed by humans, agents, and protocols. In Research
Network, the loop becomes a publishable asset: a paper, skill, workflow, evidence ledger, Walrus
snapshot, and Sui registry object moving together.

\bibliographystyle{plain}
\bibliography{references}

\end{document}
